\section{Introduction}

Bitcoin~\cite{nakamoto2008} is the most important cryptocurrency in terms of
market capitalisation. It is, however, lacking in many usability and
scalability aspects: it is costly, has low bandwidth (only a few dozen
transactions per second), and high latency. Many protocols strive to alleviate
those shortcomings.

\itparagraph{Payment Channel Networks}
Payment channels enable a pair of users to quickly and cheaply perform
arbitrarily many payments by relying on the blockchain only to resolve
disputes. Both agents pool together funds and then iteratively update their
balance. The scripts involved and the protocol itself ensure that no honest
agent loses money, by relying on a revocation mechanism: any party closing a
channel using an outdated favourable state can be punished, consequently losing
all funds invested in the channel.

Payment channels are, however, constrained to pairs of users. Two unconnected
users still need to establish a channel in order to perform a payment, which is
often not cost effective. Thus arose the concept of payment channel networks:
one transitively propagates funds across payment channels to perform a payment
between two users that are part of the same network but not connected directly.
The most popular such protocol is perhaps the Lightning
Network~\cite{poon2016lightning}, with about \$170M locked on the
network~\cite{oneml}.

Due to the sheer complexity of its specification~\cite{bolts}, it suffers from
many attacks~\cite{aumayr2021generalized, malavolta2019anonymous,
tikhomirov2020quantitative}. Despite promising ongoing
efforts~\cite{litos2026formally, grundmann2026security} it does not yet feature
a complete and unbounded verified formalisation. This paper attempts to
alleviate that issue.

\itparagraph{Protocol Verification}
Designing and analysing cryptographic protocols is a notoriously difficult and
error-prone task, as already illustrated by the various attacks and tentative
fixes on the Lightning Network. Formal methods provide hope for more robust
verification, and for automation.

Protocol verification is split between two main classes of models: the
computational model, inherited from cryptography, where attackers and
participants are polynomial-time algorithms and results are typically
probabilistic; and the symbolic or Dolev--Yao
model~\cite{dolev1983security}, where the attacker has infinite computational
power but is limited to a term algebra. The symbolic model abstracts away from
cryptography and focuses on the logic of a protocol, where much of the attacker
activity is typically found.

Although tools exist for the computational model ---
\mcryptoverif~\cite{blanchet2006computationally},
\measycrypt~\cite{barthe2011computer},
\msquirrel~\cite{baelde2021interactive},
\mcryptovampire~\cite{jeanteur2024cryptovampire} --- they are brittle, hard to
use, or narrowly scoped. Symbolic verification, by contrast, enjoys successful
automated verifiers in \mproverif~\cite{blanchet2001efficient} and
\mtamarin~\cite{meier2013tamarin}. \mtamarin\ is a multiset rewriting-based
automated protocol verifier with a pedigree of widely used protocols verified,
including WireGuard~\cite{lipp2019mechanised} and
TLS~1.3~\cite{cremers2017comprehensive}.

\itparagraph{Contributions}
We provide a modelisation of the Lightning Network in \mtamarin\ featuring
unbounded channels and routing on unbounded paths. Specifically:

\begin{itemize}
  \item \textbf{An integrated model.} We mechanise the two-party channel
  lifecycle --- including revocation and punishment --- together with multi-hop
  HTLC routing, per-hop fees, and staggered CLTV timing, in one symbolic model
  under a full Dolev--Yao adversary. To our knowledge this combination has not
  previously been machine-checked. The model is modular: five theories
  establishing 61 distinct properties across 87 lemma checks.

  \item \textbf{A linear single-spend discipline.} We identify a modelling
  pitfall in which a one-shot on-chain output is encoded as a persistent fact,
  permitting the corresponding rule to fire without bound while every safety
  lemma continues to hold. We found two independent instances in our own model
  and repaired both uniformly. The lesson generalises to any
  multiset-rewriting model of a blockchain protocol.

  \item \textbf{Quantified economic properties.} We prove end-to-end value
  conservation with strictly positive per-hop fees, and use it to pin the
  wormhole attack's theft to exactly the sum of the bypassed fees rather than
  to unspecified ``loss''.

  \item \textbf{Empirical assumption minimality.} We show by counterexample
  that the intermediary claim obligation is not implied by the timelock
  ordering constraints, and that the early-timeout race is reachable precisely
  when its ordering assumption is dropped.

  \item \textbf{Unbounded path length.} We lift the safety results to paths of
  arbitrary length by abstracting the routing layer, and give the refinement
  argument that justifies the transfer, along with an explicit account of what
  the abstraction costs.
\end{itemize}

We do \emph{not} model payment privacy or unlinkability, which are
confidentiality properties requiring onion routing and
observational-equivalence proofs. We state this as an explicit exclusion rather
than an oversight.
